% Generated by tools/build.py -- do not edit; edit data/software-engineering/*.txt
% No preamble here: this file is \input by main.tex.
\chapter{Unit 6: Software Engineering}

\section{Software Process and Prescriptive Process Models}

\begin{question}{H}
Match List I with List II.\\
List I (Process model): A. Waterfall model  B. Spiral model  C. Prototyping model  D. Incremental model\\
List II (Characteristic): I. Risk-driven; each loop includes risk analysis  II. Linear sequential phases; suits well-understood, stable requirements  III. Delivers the product in increments, the first usually being the core product  IV. Builds a quick working model to clarify unclear requirements
\optsshort{A-II, B-I, C-IV, D-III}{A-I, B-II, C-IV, D-III}{A-II, B-IV, C-I, D-III}{A-II, B-I, C-III, D-IV}
\end{question}

\begin{question}{H}
The spiral model was proposed by
\optsshort{Watts Humphrey}{Barry Boehm}{Ivar Jacobson}{Winston Royce}
\end{question}

\begin{question}{H}
The spiral model is considered a meta-model because
\optslong{it has no phases}{it uses metadata}{it can incorporate other process models such as waterfall and prototyping}{it is used only for maintenance}
\end{question}

\begin{question}{H}
Arrange the phases of the Unified Process in order.\\
A. Construction  B. Elaboration  C. Inception  D. Transition
\optsshort{C, B, D, A}{C, B, A, D}{B, C, A, D}{C, A, B, D}
\end{question}

\begin{question}{H}
Given below are two statements.\\
Statement I: The Unified Process is use-case driven, architecture-centric, iterative and incremental.\\
Statement II: The Unified Process was developed by Jacobson, Booch and Rumbaugh.
\optslong{Both Statement I and Statement II are true}{Both Statement I and Statement II are false}{Statement I is true but Statement II is false}{Statement I is false but Statement II is true}
\end{question}

\begin{question}{H}
Pressman's generic process framework consists of the five framework activities
\optslong{requirements, design, coding, testing, maintenance}{analysis, design, implementation, integration, operation}{communication, planning, modeling, construction, deployment}{inception, elaboration, construction, transition, production}
\end{question}

\begin{question}{H}
Which of the following are umbrella activities in Pressman's process framework?\\
A. Software project tracking and control  B. Risk management  C. Software quality assurance  D. Software configuration management  E. Construction\\
Choose the correct answer from the options given below:
\optsshort{A, B and E only}{C, D and E only}{A, B, C, D and E}{A, B, C and D only}
\end{question}

\begin{question}{H}
The V-model differs from the waterfall model in that it
\optslong{uses prototypes}{is iterative}{has no testing}{explicitly pairs each development phase with a corresponding testing phase}
\end{question}

\begin{question}{H}
Given below are two statements: one is labelled as Assertion A and the other is labelled as Reason R.\\
Assertion A: The waterfall model is unsuitable when requirements are likely to change.\\
Reason R: In the waterfall model, each phase must be completed before the next begins, and going back to earlier phases is costly.
\optslong{Both A and R are true and R is the correct explanation of A}{Both A and R are true but R is NOT the correct explanation of A}{A is true but R is false}{A is false but R is true}
\end{question}

\begin{question}{M}
The Rapid Application Development (RAD) model is suitable when
\optslong{requirements are unclear and the system is very large}{requirements are well understood, the project scope is constrained, and the system can be modularised, with a typical cycle of 60--90 days}{technical risks are very high}{there is no customer involvement}
\end{question}

\begin{question}{M}
Given below are two statements.\\
Statement I: In throwaway prototyping, the prototype is discarded after requirements are understood.\\
Statement II: In evolutionary prototyping, the prototype is refined into the final system.
\optslong{Both Statement I and Statement II are true}{Both Statement I and Statement II are false}{Statement I is true but Statement II is false}{Statement I is false but Statement II is true}
\end{question}

\begin{question}{M}
In the concurrent development model, activities are represented as
\optslong{a spiral}{a single loop}{a series of states (such as under development, awaiting changes, done) with events triggering transitions}{a linear sequence}
\end{question}

\begin{question}{M}
Component-based development emphasises
\optslong{building systems by integrating pre-packaged, reusable software components}{formal mathematical specification}{pair programming}{writing all code from scratch}
\end{question}

\begin{question}{M}
Aspect-oriented software development addresses
\optslong{database design}{hardware design}{only user interfaces}{crosscutting concerns such as logging and security that span many modules}
\end{question}

\begin{question}{H}
Given below are two statements.\\
Statement I: The formal methods model uses mathematical notation to specify, develop and verify software.\\
Statement II: The cleanroom approach to software engineering is a variation of the formal methods model.
\optslong{Both Statement I and Statement II are true}{Both Statement I and Statement II are false}{Statement I is true but Statement II is false}{Statement I is false but Statement II is true}
\end{question}

\begin{question}{M}
Which of the following is a formal specification language?
\optsshort{SQL}{Z notation}{HTML}{UML}
\end{question}

\begin{question}{M}
A process pattern describes
\optslong{a test case}{a user interface layout}{a process-related problem encountered during software engineering, its environment, and a proven solution}{a design pattern for classes}
\end{question}

\begin{question}{M}
The waterfall model was first described by
\optsshort{Barry Boehm}{Frederick Brooks}{Winston Royce}{Tom DeMarco}
\end{question}

\begin{question}{M}
In the incremental model, the first increment is usually
\optslong{a throwaway prototype}{the core product, addressing basic requirements}{the complete system}{documentation only}
\end{question}

\begin{question}{H}
Match List I with List II.\\
List I (Spiral model quadrant activity): A. Determine objectives, alternatives and constraints  B. Evaluate alternatives; identify and resolve risks  C. Develop and verify the next-level product  D. Plan the next phases\\
List II (Order in a loop): I. Third  II. First  III. Fourth  IV. Second
\optsshort{A-II, B-IV, C-I, D-III}{A-II, B-I, C-IV, D-III}{A-II, B-IV, C-III, D-I}{A-IV, B-II, C-I, D-III}
\end{question}

\begin{question}{M}
Given below are two statements: one is labelled as Assertion A and the other is labelled as Reason R.\\
Assertion A: The prototyping model is useful when the customer cannot state detailed requirements.\\
Reason R: The spiral model was proposed by Barry Boehm.
\optslong{Both A and R are true and R is the correct explanation of A}{Both A and R are true but R is NOT the correct explanation of A}{A is true but R is false}{A is false but R is true}
\end{question}

\begin{question}{M}
Which process model is best suited to a large, mission-critical project with high technical risk?
\optsshort{Big bang}{Waterfall}{Spiral}{Code-and-fix}
\end{question}

\begin{question}{M}
The Personal Software Process (PSP) and Team Software Process (TSP) were developed by
\optsshort{Watts Humphrey}{Barry Boehm}{Grady Booch}{Kent Beck}
\end{question}

\begin{question}{M}
Given below are two statements.\\
Statement I: A software process is the same as software engineering.\\
Statement II: A software process defines a framework of activities, actions and tasks needed to build high-quality software.
\optslong{Both Statement I and Statement II are true}{Both Statement I and Statement II are false}{Statement I is true but Statement II is false}{Statement I is false but Statement II is true}
\end{question}

\begin{question}{M}
Software engineering, according to IEEE, is
\optslong{managing hardware}{testing software only}{writing programs quickly}{the application of a systematic, disciplined, quantifiable approach to the development, operation and maintenance of software}
\end{question}

\begin{question}{M}
Arrange the layers of software engineering as described by Pressman, from the bottom (foundation) to the top.\\
A. Methods  B. Tools  C. A quality focus  D. Process
\optsshort{C, D, A, B}{C, A, D, B}{D, C, A, B}{B, A, D, C}
\end{question}

\begin{question}{M}
Given below are two statements.\\
Statement I: The code-and-fix approach works well for large projects with many developers.\\
Statement II: The big-bang model starts with detailed planning and formal requirements.
\optslong{Both Statement I and Statement II are true}{Both Statement I and Statement II are false}{Statement I is true but Statement II is false}{Statement I is false but Statement II is true}
\end{question}

\begin{question}{L}
A software product that is sold to many customers in the open market is called a
\optsshort{legacy system}{generic product}{bespoke (custom) product}{prototype}
\end{question}

\begin{question}{M}
Which of the following is a characteristic of software that distinguishes it from hardware?
\optslong{Software has a bathtub failure curve identical to hardware}{Software is always mass-produced}{Software wears out with use}{Software is developed (engineered), not manufactured, and does not wear out in the classical sense}
\end{question}

\begin{question}{M}
Given below are two statements: one is labelled as Assertion A and the other is labelled as Reason R.\\
Assertion A: The cost of fixing a defect increases the later it is found in the life cycle.\\
Reason R: Defects found late require changes to more work products, such as design, code, tests and documentation.
\optslong{Both A and R are true and R is the correct explanation of A}{Both A and R are true but R is NOT the correct explanation of A}{A is true but R is false}{A is false but R is true}
\end{question}

\begin{question}{M}
In the Unified Process, most use cases are captured and the architecture baseline is established during
\optsshort{inception}{transition}{construction}{elaboration}
\end{question}

\begin{answers}
\ans{1}{1}{Standard prescriptive process models in Pressman.}
\ans{2}{2}{Boehm's spiral model (1986--88) combines iterative prototyping with systematic risk analysis.}
\ans{3}{3}{Each cycle can follow whichever model best reduces the current risks.}
\ans{4}{2}{Pressman also lists a production phase after transition.}
\ans{5}{1}{The Rational Unified Process (RUP) is its best-known refinement.}
\ans{6}{3}{Umbrella activities (such as risk management and SQA) apply across all of them.}
\ans{7}{4}{Construction is a framework activity, not an umbrella activity.}
\ans{8}{4}{For example, unit testing verifies detailed design and acceptance testing validates requirements.}
\ans{9}{1}{Real projects rarely follow a strictly sequential flow.}
\ans{10}{2}{RAD relies on component reuse and multiple teams working in parallel.}
\ans{11}{1}{Throwaway prototypes risk being kept for convenience, which leads to poor quality.}
\ans{12}{3}{It is useful when several activities are going on at the same time.}
\ans{13}{1}{Steps include researching available components, integration issues and architecture design.}
\ans{14}{4}{Aspects modularise concerns that would otherwise be scattered across components.}
\ans{15}{1}{Formal methods are costly and require training, but they can eliminate ambiguity and many defects.}
\ans{16}{2}{VDM and B are other formal specification languages.}
\ans{17}{3}{Types of process pattern are stage, task and phase patterns.}
\ans{18}{3}{Royce's 1970 paper actually recommended iteration and feedback.}
\ans{19}{2}{Later increments add supplementary features based on feedback.}
\ans{20}{1}{Each loop of the spiral passes through these four quadrants.}
\ans{21}{2}{Both are true, but R is unrelated; prototyping helps because the customer and developer can see and refine a working model.}
\ans{22}{3}{Its explicit risk analysis in each cycle suits high-risk projects.}
\ans{23}{1}{They emphasise personal measurement and disciplined team processes.}
\ans{24}{4}{Software engineering encompasses the process, methods and tools.}
\ans{25}{4}{Pressman describes it as a layered technology: quality focus, process, methods and tools.}
\ans{26}{1}{The bedrock is an organisational commitment to quality.}
\ans{27}{2}{Code-and-fix and big-bang approaches suit only tiny projects and involve little planning.}
\ans{28}{2}{Bespoke software is developed for a particular customer.}
\ans{29}{4}{Software's failure rate rises over time because of side effects of changes, giving an idealised and an actual curve.}
\ans{30}{1}{This is the motivation for early reviews and testing.}
\ans{31}{4}{Inception defines the business case; construction builds the software.}
\end{answers}

\section{Agile Process Models (Agile Manifesto, XP, Scrum, ASD, DSDM, FDD, Crystal, Web Engineering)}

\begin{question}{H}
Which of the following are values stated in the Agile Manifesto?\\
A. Individuals and interactions over processes and tools  B. Working software over comprehensive documentation  C. Customer collaboration over contract negotiation  D. Following a plan over responding to change\\
Choose the correct answer from the options given below:
\optsshort{A and D only}{B, C and D only}{A, B and C only}{A, B, C and D}
\end{question}

\begin{question}{H}
Match List I with List II.\\
List I (Agile method): A. Extreme Programming (XP)  B. Scrum  C. Adaptive Software Development (ASD)  D. Feature Driven Development (FDD)\\
List II (Key characteristic): I. Speculation, collaboration and learning cycles  II. Pair programming, test-first development and refactoring  III. Client-valued features deliverable in two weeks or less  IV. Sprints, backlogs and daily stand-up meetings
\optsshort{A-II, B-I, C-IV, D-III}{A-II, B-IV, C-I, D-III}{A-IV, B-II, C-I, D-III}{A-II, B-IV, C-III, D-I}
\end{question}

\begin{question}{H}
Match List I with List II.\\
List I (Scrum element): A. Product backlog  B. Sprint  C. Daily Scrum  D. Sprint retrospective\\
List II (Description): I. Short daily meeting (about 15 minutes) to synchronise the team  II. Prioritised list of features required by the customer  III. Time-boxed iteration, typically 2--4 weeks, producing a potentially shippable increment  IV. Meeting where the team reflects on how to improve its process
\optsshort{A-II, B-III, C-IV, D-I}{A-III, B-II, C-I, D-IV}{A-II, B-I, C-III, D-IV}{A-II, B-III, C-I, D-IV}
\end{question}

\begin{question}{H}
The three roles defined in Scrum are
\optslong{Product Owner, Scrum Master and Development Team}{customer, analyst and coder}{architect, designer and programmer}{project manager, tester and developer}
\end{question}

\begin{question}{H}
Which of the following practices belong to Extreme Programming?\\
A. Pair programming  B. Test-first development  C. Refactoring  D. Continuous integration  E. Comprehensive up-front design documentation\\
Choose the correct answer from the options given below:
\optsshort{A, B and E only}{A, B, C and D only}{C, D and E only}{A, B, C, D and E}
\end{question}

\begin{question}{H}
Given below are two statements.\\
Statement I: In XP, a spike solution is an operational prototype built to reduce risk on a difficult design problem.\\
Statement II: In XP, user stories are written by the customer to describe required features.
\optslong{Both Statement I and Statement II are true}{Both Statement I and Statement II are false}{Statement I is true but Statement II is false}{Statement I is false but Statement II is true}
\end{question}

\begin{question}{M}
The five values of Extreme Programming are
\optslong{planning, design, coding, testing and release}{speed, cost, quality, scope and time}{communication, simplicity, feedback, courage and respect}{individuals, software, collaboration, change and plans}
\end{question}

\begin{question}{M}
CRC cards, used in XP and object-oriented analysis, stand for
\optslong{Code-Review-Check}{Class-Relation-Constraint}{Class-Responsibility-Collaborator}{Customer-Requirement-Change}
\end{question}

\begin{question}{H}
DSDM (Dynamic Systems Development Method) is based on the observation that
\optslong{no iteration is needed}{documentation is more important than code}{100\% of the application must be delivered at once}{80\% of an application can often be delivered in 20\% of the time it would take to deliver 100\%}
\end{question}

\begin{question}{M}
In MoSCoW prioritisation, used in DSDM, the letters stand for
\optslong{Measure, Schedule, Control, Watch}{Modules, Systems, Components, Widgets}{Model, Specify, Code, Write}{Must have, Should have, Could have, Won't have this time}
\end{question}

\begin{question}{M}
In FDD, a feature is expressed using the template
\optslong{<action> the <result> <by | for | of | to> a(n) <object>}{<noun> <verb> <noun>}{as a <role> I want <goal>}{given <context> when <event> then <outcome>}
\end{question}

\begin{question}{M}
Arrange the five processes of Feature Driven Development in order.\\
A. Build a features list  B. Design by feature  C. Develop an overall model  D. Build by feature  E. Plan by feature
\optsshort{C, A, E, B, D}{C, E, A, B, D}{A, C, E, B, D}{C, A, B, E, D}
\end{question}

\begin{question}{M}
The Crystal family of agile methods, proposed by Alistair Cockburn, varies its methods according to
\optslong{the hardware used}{the programming language}{the customer's country}{project size (team size) and criticality, indicated by colours such as Clear, Yellow and Orange}
\end{question}

\begin{question}{M}
Adaptive Software Development (ASD), proposed by Jim Highsmith, has a life cycle of
\optslong{planning, design and coding}{speculation, collaboration and learning}{analysis, design and testing}{inception, elaboration and construction}
\end{question}

\begin{question}{H}
Given below are two statements: one is labelled as Assertion A and the other is labelled as Reason R.\\
Assertion A: Agile methods deliver working software frequently in short iterations.\\
Reason R: Agile methods assume that requirements are fully known and fixed at the start of the project.
\optslong{Both A and R are true and R is the correct explanation of A}{Both A and R are true but R is NOT the correct explanation of A}{A is true but R is false}{A is false but R is true}
\end{question}

\begin{question}{M}
A burndown chart in Scrum shows
\optslong{the number of defects found per day}{the team's salaries}{the architecture of the system}{the amount of work remaining in a sprint or release plotted against time}
\end{question}

\begin{question}{M}
Given below are two statements.\\
Statement I: In Scrum, the Product Owner is responsible for maximising the value of the product and managing the product backlog.\\
Statement II: In Scrum, the Scrum Master assigns tasks to individual developers.
\optslong{Both Statement I and Statement II are true}{Both Statement I and Statement II are false}{Statement I is true but Statement II is false}{Statement I is false but Statement II is true}
\end{question}

\begin{question}{M}
Test-driven development (TDD) follows the cycle
\optslong{test only at the end of the project}{write a failing test, write code to make it pass, then refactor}{no testing}{write code, then write tests later}
\end{question}

\begin{question}{M}
Which of the following is NOT one of the twelve principles behind the Agile Manifesto?
\optslong{Complete all documentation before coding begins}{Business people and developers must work together daily}{Maintain a sustainable pace indefinitely}{Deliver working software frequently}
\end{question}

\begin{question}{M}
Web applications (WebApps) differ from traditional software in that they are typically
\optslong{network intensive, concurrent, subject to unpredictable load, content driven and continuously evolving}{free of security concerns}{never updated}{used by one user at a time}
\end{question}

\begin{question}{M}
Kanban, used in lean and agile development, emphasises
\optslong{detailed up-front planning}{visualising work on a board and limiting work in progress}{fixed-length sprints}{pair programming}
\end{question}

\begin{question}{M}
Given below are two statements: one is labelled as Assertion A and the other is labelled as Reason R.\\
Assertion A: Pair programming can improve code quality.\\
Reason R: In pair programming, two developers work independently on separate modules.
\optslong{Both A and R are true and R is the correct explanation of A}{Both A and R are true but R is NOT the correct explanation of A}{A is true but R is false}{A is false but R is true}
\end{question}

\begin{question}{M}
In XP, refactoring means
\optslong{improving the internal structure of code without changing its external behaviour}{adding new features}{rewriting the program in another language}{removing tests}
\end{question}

\begin{question}{L}
The Agile Manifesto was published in
\optsshort{2001}{2011}{1995}{1991}
\end{question}

\begin{question}{M}
Given below are two statements.\\
Statement I: Agile methods prohibit all forms of documentation.\\
Statement II: Agile methods are well suited to small and medium teams working with customers who can collaborate closely.
\optslong{Both Statement I and Statement II are true}{Both Statement I and Statement II are false}{Statement I is true but Statement II is false}{Statement I is false but Statement II is true}
\end{question}

\begin{question}{M}
In XP, the practice of having all code integrated and tested several times a day is called
\optsshort{pair programming}{small releases}{continuous integration}{collective ownership}
\end{question}

\begin{answers}
\ans{1}{3}{The fourth value is responding to change over following a plan.}
\ans{2}{2}{Characteristic practices of common agile methods.}
\ans{3}{4}{Scrum artifacts and events.}
\ans{4}{1}{The Scrum Master facilitates the process and removes impediments.}
\ans{5}{2}{XP favours simple design and minimal documentation.}
\ans{6}{1}{Stories are prioritised and estimated by the team.}
\ans{7}{3}{XP practices are derived from these values.}
\ans{8}{3}{Each card lists a class, its responsibilities and the classes it collaborates with.}
\ans{9}{4}{It uses time-boxing and MoSCoW prioritisation.}
\ans{10}{4}{It helps decide what fits into a time box.}
\ans{11}{1}{For example, ``Add the product to a shopping cart''.}
\ans{12}{1}{The first three processes are done once; the last two iterate per feature.}
\ans{13}{4}{More people and higher criticality call for heavier methods.}
\ans{14}{2}{It emphasises self-organising teams and learning.}
\ans{15}{3}{Agile methods assume requirements will change and welcome change even late in development.}
\ans{16}{4}{It helps the team see whether it is on track to complete the sprint.}
\ans{17}{3}{The Development Team is self-organising and chooses how to do its work; the Scrum Master is a facilitator.}
\ans{18}{2}{This is often summarised as red, green, refactor.}
\ans{19}{1}{Agile values working software over comprehensive documentation.}
\ans{20}{1}{These attributes shape the web engineering (WebE) process.}
\ans{21}{2}{It comes from the Toyota Production System.}
\ans{22}{3}{In pair programming two developers share one workstation: one writes code while the other continuously reviews it.}
\ans{23}{1}{Refactoring keeps the design simple as the code evolves.}
\ans{24}{1}{It was written by seventeen practitioners at Snowbird, Utah.}
\ans{25}{4}{Agile methods value working software more than documentation but still produce documentation that adds value.}
\ans{26}{3}{It detects integration problems early.}
\end{answers}

\section{Software Requirements (functional and non-functional requirements, elicitation, use cases, analysis modelling, SRS)}

\begin{question}{H}
Match List I with List II.\\
List I (Requirement): A. The system shall allow a user to reset the password  B. The system shall respond to any query within 2 seconds  C. The system shall run on Linux and Windows  D. The system shall encrypt all stored passwords\\
List II (Type): I. Portability requirement  II. Functional requirement  III. Security requirement  IV. Performance requirement
\optsshort{A-II, B-I, C-IV, D-III}{A-IV, B-II, C-I, D-III}{A-II, B-IV, C-I, D-III}{A-II, B-IV, C-III, D-I}
\end{question}

\begin{question}{H}
Given below are two statements.\\
Statement I: Functional requirements describe the services the system should provide.\\
Statement II: Non-functional requirements describe constraints such as performance, reliability and usability.
\optslong{Both Statement I and Statement II are true}{Both Statement I and Statement II are false}{Statement I is true but Statement II is false}{Statement I is false but Statement II is true}
\end{question}

\begin{question}{H}
Arrange the requirements engineering tasks described by Pressman in order.\\
A. Elicitation  B. Inception  C. Negotiation  D. Elaboration  E. Specification
\optsshort{B, D, A, C, E}{B, A, C, D, E}{A, B, D, C, E}{B, A, D, C, E}
\end{question}

\begin{question}{H}
Which of the following are characteristics of a good SRS according to IEEE 830?\\
A. Correct  B. Unambiguous  C. Complete  D. Verifiable  E. Describes how the system will be implemented\\
Choose the correct answer from the options given below:
\optsshort{A, B, C, D and E}{C, D and E only}{A, B and E only}{A, B, C and D only}
\end{question}

\begin{question}{H}
Given below are two statements: one is labelled as Assertion A and the other is labelled as Reason R.\\
Assertion A: Every requirement in an SRS should be verifiable.\\
Reason R: A requirement such as ``the system should be user-friendly'' cannot be tested objectively.
\optslong{Both A and R are true and R is the correct explanation of A}{Both A and R are true but R is NOT the correct explanation of A}{A is true but R is false}{A is false but R is true}
\end{question}

\begin{question}{H}
In Quality Function Deployment (QFD), requirements that the customer does not explicitly state but whose absence causes dissatisfaction are called
\optsshort{normal requirements}{expected requirements}{exciting requirements}{functional requirements}
\end{question}

\begin{question}{H}
Match List I with List II.\\
List I (Analysis model element): A. Scenario-based  B. Class-based  C. Behavioral  D. Flow-oriented\\
List II (Example): I. State diagrams and sequence diagrams  II. Use cases and activity diagrams  III. Data flow diagrams and the data dictionary  IV. Class diagrams and CRC models
\optsshort{A-IV, B-II, C-I, D-III}{A-II, B-I, C-IV, D-III}{A-II, B-IV, C-I, D-III}{A-II, B-IV, C-III, D-I}
\end{question}

\begin{question}{H}
A level-0 data flow diagram (context diagram) shows
\optslong{only data stores}{the internal processes in detail}{the control flow of the program}{the entire system as a single process (bubble) with its external entities and data flows}
\end{question}

\begin{question}{M}
Given below are two statements.\\
Statement I: A data flow diagram shows the sequence and control logic of processing.\\
Statement II: When a DFD bubble is refined, the inputs and outputs of the refined diagram must match those of the parent bubble (balancing).
\optslong{Both Statement I and Statement II are true}{Both Statement I and Statement II are false}{Statement I is true but Statement II is false}{Statement I is false but Statement II is true}
\end{question}

\begin{question}{M}
In a use case diagram, the relationship in which a base use case optionally includes extra behaviour under certain conditions is called
\optsshort{association}{generalization}{extend}{include}
\end{question}

\begin{question}{M}
An actor in a use case diagram is
\optslong{a class in the system}{a database table}{a role played by a person or external system that interacts with the system}{a test case}
\end{question}

\begin{question}{H}
Which of the following is a technique for requirements elicitation?\\
A. Interviews  B. Collaborative requirements gathering (JAD)  C. Questionnaires  D. Observation (ethnography)  E. Code inspection\\
Choose the correct answer from the options given below:
\optsshort{C, D and E only}{A, B, C and D only}{A, B, C, D and E}{A, B and E only}
\end{question}

\begin{question}{M}
A requirements traceability matrix is used to
\optslong{draw DFDs}{estimate cost}{link each requirement to its sources, design elements, code and test cases}{schedule the project}
\end{question}

\begin{question}{M}
The data dictionary in structured analysis contains
\optslong{test results}{the source code}{definitions of all data elements, data flows and data stores used in the DFDs}{the project schedule}
\end{question}

\begin{question}{M}
Given below are two statements: one is labelled as Assertion A and the other is labelled as Reason R.\\
Assertion A: Requirements validation checks that the requirements actually define the system the customer wants.\\
Reason R: Requirements errors found after delivery are much cheaper to fix than errors found during the requirements phase.
\optslong{Both A and R are true and R is the correct explanation of A}{Both A and R are true but R is NOT the correct explanation of A}{A is true but R is false}{A is false but R is true}
\end{question}

\begin{question}{M}
Which of the following is NOT a requirements validation technique?
\optsshort{Test-case generation}{Prototyping}{Regression testing}{Requirements reviews}
\end{question}

\begin{question}{M}
Domain requirements are requirements that
\optslong{are always non-functional}{are optional}{apply only to the user interface}{come from the application domain of the system and reflect its characteristics}
\end{question}

\begin{question}{M}
In a swimlane diagram (a variation of the activity diagram), each lane represents
\optslong{a time period}{a use case}{an actor or analysis class responsible for the activities in it}{a database table}
\end{question}

\begin{question}{M}
Given below are two statements.\\
Statement I: Once the SRS is approved, requirements never change in practice.\\
Statement II: Requirements management is the process of managing changes to requirements throughout the project.
\optslong{Both Statement I and Statement II are true}{Both Statement I and Statement II are false}{Statement I is true but Statement II is false}{Statement I is false but Statement II is true}
\end{question}

\begin{question}{M}
Ambiguity in requirements can be reduced most effectively by
\optslong{starting coding early}{avoiding customer contact}{using precise, measurable language, glossaries and, where needed, formal specifications}{writing longer documents}
\end{question}

\begin{question}{M}
The IEEE standard for the Software Requirements Specification is
\optsshort{IEEE 1012}{IEEE 730}{IEEE 829}{IEEE 830}
\end{question}

\begin{question}{M}
The Kano model classifies customer requirements into
\optslong{hardware and software}{functional and non-functional only}{basic (must-be), performance and excitement (delighter) requirements}{internal and external}
\end{question}

\begin{question}{M}
Given below are two statements: one is labelled as Assertion A and the other is labelled as Reason R.\\
Assertion A: Use cases are useful for capturing functional requirements.\\
Reason R: UML is standardised by the Object Management Group.
\optslong{Both A and R are true and R is the correct explanation of A}{Both A and R are true but R is NOT the correct explanation of A}{A is true but R is false}{A is false but R is true}
\end{question}

\begin{question}{M}
Feasibility study in requirements engineering assesses
\optslong{the cyclomatic complexity}{the coupling of modules}{the number of test cases}{whether the system is technically, economically and operationally feasible}
\end{question}

\begin{question}{L}
SRS stands for
\optslong{System Resource Schedule}{Software Release Summary}{Software Requirements Specification}{Structured Requirement Syntax}
\end{question}

\begin{question}{M}
Given below are two statements.\\
Statement I: An entity-relationship diagram can be used in data modelling during requirements analysis.\\
Statement II: State diagrams model the behaviour of a system in response to external events.
\optslong{Both Statement I and Statement II are true}{Both Statement I and Statement II are false}{Statement I is true but Statement II is false}{Statement I is false but Statement II is true}
\end{question}

\begin{question}{M}
Requirements negotiation is needed because
\optslong{testing is expensive}{requirements are always consistent}{customers and stakeholders often have conflicting requirements and limited resources}{the developers write the requirements}
\end{question}

\begin{answers}
\ans{1}{3}{Functional requirements describe what the system does; the others are non-functional (quality) requirements.}
\ans{2}{1}{Non-functional requirements often apply to the system as a whole.}
\ans{3}{4}{Validation and requirements management follow specification.}
\ans{4}{4}{An SRS states what the system must do, not how it will be designed.}
\ans{5}{1}{Requirements should be stated in measurable terms, for example ``a new user completes registration within 3 minutes''.}
\ans{6}{2}{Normal requirements are stated explicitly; exciting requirements go beyond expectations.}
\ans{7}{3}{Elements of the requirements (analysis) model in Pressman.}
\ans{8}{4}{Lower levels refine the single bubble into sub-processes.}
\ans{9}{4}{A DFD shows the flow of data, not control or sequencing.}
\ans{10}{3}{An include relationship is mandatory: the base use case always uses the included behaviour.}
\ans{11}{3}{Primary actors initiate use cases; secondary actors support the system.}
\ans{12}{2}{Code inspection is a verification technique, not elicitation.}
\ans{13}{3}{It helps assess the impact of changes and confirm that every requirement is tested.}
\ans{14}{3}{It gives each name a precise meaning.}
\ans{15}{3}{Errors found later are far more expensive, which is why validation is done early.}
\ans{16}{3}{Regression testing is done after changes to code.}
\ans{17}{4}{For example, rules and regulations of banking for a banking system.}
\ans{18}{3}{It shows who is responsible for each activity.}
\ans{19}{4}{Requirements change frequently; change control and traceability are needed.}
\ans{20}{3}{Each requirement should have only one interpretation.}
\ans{21}{4}{IEEE 829 is for test documentation and IEEE 730 for SQA plans.}
\ans{22}{3}{It is related to the normal, expected and exciting requirements of QFD.}
\ans{23}{2}{Both are true, but use cases help because they describe actor-system interactions from the user's point of view.}
\ans{24}{4}{It decides whether the project is worth pursuing.}
\ans{25}{3}{It is the official statement of what the developers should implement.}
\ans{26}{1}{Both are parts of the analysis model.}
\ans{27}{3}{The win-win approach seeks requirements acceptable to all stakeholders.}
\end{answers}

\section{Software Design (design concepts, cohesion and coupling, architectural, data, user-interface and component-level design, design patterns)}

\begin{question}{H}
Arrange the following types of cohesion from the lowest (worst) to the highest (best).\\
A. Functional  B. Coincidental  C. Temporal  D. Sequential  E. Logical
\optsshort{B, C, E, D, A}{B, E, C, D, A}{B, E, D, C, A}{E, B, C, D, A}
\end{question}

\begin{question}{H}
Arrange the following types of coupling from the worst (tightest) to the best (loosest).\\
A. Data coupling  B. Content coupling  C. Control coupling  D. Common coupling  E. Stamp coupling
\optsshort{B, D, E, C, A}{B, D, C, E, A}{D, B, C, E, A}{B, C, D, E, A}
\end{question}

\begin{question}{H}
Match List I with List II.\\
List I (Coupling type): A. Content coupling  B. Common coupling  C. Control coupling  D. Data coupling\\
List II (Description): I. Modules share global data  II. One module passes a flag that controls the logic of another  III. Modules communicate only through parameters that are elementary data items  IV. One module directly modifies or depends on the internal workings of another
\optsshort{A-I, B-IV, C-II, D-III}{A-IV, B-II, C-I, D-III}{A-IV, B-I, C-II, D-III}{A-IV, B-I, C-III, D-II}
\end{question}

\begin{question}{H}
Match List I with List II.\\
List I (Cohesion type): A. Logical cohesion  B. Temporal cohesion  C. Communicational cohesion  D. Functional cohesion\\
List II (Description): I. All elements contribute to a single well-defined task  II. Elements perform logically similar tasks, selected by a control flag  III. Elements operate on the same input or output data  IV. Elements are grouped because they execute at the same time, such as initialisation
\optsshort{A-II, B-III, C-IV, D-I}{A-IV, B-II, C-III, D-I}{A-II, B-IV, C-III, D-I}{A-II, B-IV, C-I, D-III}
\end{question}

\begin{question}{H}
Given below are two statements.\\
Statement I: A good design aims for high cohesion and low coupling.\\
Statement II: Functional independence is achieved by modules with a single purpose and limited interaction with other modules.
\optslong{Both Statement I and Statement II are true}{Both Statement I and Statement II are false}{Statement I is true but Statement II is false}{Statement I is false but Statement II is true}
\end{question}

\begin{question}{H}
The concept of information hiding, proposed by David Parnas, states that
\optslong{modules should hide their design decisions and data from other modules, exposing only what is necessary through interfaces}{code should be obfuscated}{all data should be global}{programs should not have comments}
\end{question}

\begin{question}{H}
Given below are two statements: one is labelled as Assertion A and the other is labelled as Reason R.\\
Assertion A: Dividing software into a very large number of tiny modules minimises total development cost.\\
Reason R: As the number of modules grows, the cost of integrating them (interfaces) also grows.
\optslong{Both A and R are true and R is the correct explanation of A}{Both A and R are true but R is NOT the correct explanation of A}{A is true but R is false}{A is false but R is true}
\end{question}

\begin{question}{H}
Match List I with List II.\\
List I (Architectural style): A. Data-centred  B. Data-flow  C. Call-and-return  D. Layered\\
List II (Example): I. Pipes and filters  II. Main program and subprograms, remote procedure call  III. A central repository or blackboard accessed by independent clients  IV. OSI protocol stack with services at each level
\optsshort{A-III, B-I, C-II, D-IV}{A-I, B-III, C-II, D-IV}{A-III, B-I, C-IV, D-II}{A-III, B-II, C-I, D-IV}
\end{question}

\begin{question}{H}
Match List I with List II.\\
List I (Design principle): A. Single responsibility principle  B. Open-closed principle  C. Liskov substitution principle  D. Dependency inversion principle\\
List II (Statement): I. Subclasses should be substitutable for their base classes  II. Depend on abstractions, not on concrete classes  III. A class should have only one reason to change  IV. A module should be open for extension but closed for modification
\optsshort{A-III, B-I, C-IV, D-II}{A-III, B-IV, C-II, D-I}{A-IV, B-III, C-I, D-II}{A-III, B-IV, C-I, D-II}
\end{question}

\begin{question}{H}
Match List I with List II.\\
List I (Design pattern): A. Singleton  B. Observer  C. Adapter  D. Factory method\\
List II (Category and purpose): I. Behavioral; notifies dependents automatically when an object changes state  II. Creational; lets subclasses decide which class to instantiate  III. Creational; ensures a class has only one instance  IV. Structural; converts the interface of a class into one that clients expect
\optsshort{A-III, B-I, C-II, D-IV}{A-III, B-IV, C-I, D-II}{A-I, B-III, C-IV, D-II}{A-III, B-I, C-IV, D-II}
\end{question}

\begin{question}{M}
The Gang of Four catalogue classifies its 23 design patterns into
\optslong{architectural, process and test patterns}{static and dynamic patterns}{creational, structural and behavioral patterns}{analysis, design and coding patterns}
\end{question}

\begin{question}{M}
The Model-View-Controller (MVC) architecture separates
\optslong{the database from the network}{hardware from software}{data and business logic (model), presentation (view) and input handling (controller)}{testing from development}
\end{question}

\begin{question}{H}
Theo Mandel's three golden rules for user-interface design are
\optslong{use only keyboard input}{use many colours, many fonts and many windows}{place the user in control, reduce the user's memory load, and make the interface consistent}{minimise testing, maximise features, use defaults}
\end{question}

\begin{question}{M}
Stepwise refinement, a top-down design strategy, was proposed by
\optsshort{Niklaus Wirth}{David Parnas}{Barry Boehm}{Edsger Dijkstra}
\end{question}

\begin{question}{M}
Given below are two statements.\\
Statement I: Abstraction allows a designer to focus on a problem at some level of generalisation without regard to low-level details.\\
Statement II: Refactoring changes the external behaviour of a program to add new features.
\optslong{Both Statement I and Statement II are true}{Both Statement I and Statement II are false}{Statement I is true but Statement II is false}{Statement I is false but Statement II is true}
\end{question}

\begin{question}{M}
In a structure chart, fan-out of a module is
\optslong{the number of modules it directly calls}{the number of parameters}{the depth of the chart}{the number of modules that call it}
\end{question}

\begin{question}{M}
Given below are two statements: one is labelled as Assertion A and the other is labelled as Reason R.\\
Assertion A: High fan-out is generally desirable in a module structure.\\
Reason R: A module with high fan-in is reused by many other modules.
\optslong{Both A and R are true and R is the correct explanation of A}{Both A and R are true but R is NOT the correct explanation of A}{A is true but R is false}{A is false but R is true}
\end{question}

\begin{question}{M}
In structured design, transform analysis is used when the data flow diagram shows
\optslong{no data flow}{a single item triggering one of many paths (a transaction centre)}{only data stores}{incoming flow, a transform centre and outgoing flow}
\end{question}

\begin{question}{M}
The facade design pattern
\optslong{defines a family of algorithms}{ensures one instance}{provides a unified, simplified interface to a set of interfaces in a subsystem}{lets an object change its behaviour when its state changes}
\end{question}

\begin{question}{M}
The interface segregation principle states that
\optslong{clients should not be forced to depend on interfaces they do not use}{interfaces must not change}{every class must implement every interface}{interfaces should be as large as possible}
\end{question}

\begin{question}{M}
Which design pattern defines a family of algorithms, encapsulates each one and makes them interchangeable?
\optsshort{Strategy}{Builder}{Composite}{Proxy}
\end{question}

\begin{question}{M}
Given below are two statements.\\
Statement I: Architectural design defines the relationships among the major structural elements of the software.\\
Statement II: Component-level design transforms structural elements into a procedural description of each component.
\optslong{Both Statement I and Statement II are true}{Both Statement I and Statement II are false}{Statement I is true but Statement II is false}{Statement I is false but Statement II is true}
\end{question}

\begin{question}{M}
Data design at the architectural level may involve
\optslong{writing test cases}{choosing loop variables}{selecting the colour scheme}{designing a data warehouse or database to support business needs}
\end{question}

\begin{question}{M}
The design concept ``separation of concerns'' means that
\optslong{data and code are always mixed}{a complex problem is easier to handle if it is subdivided into pieces that can each be solved independently}{only one person designs the system}{all concerns are placed in one module}
\end{question}

\begin{question}{M}
Given below are two statements: one is labelled as Assertion A and the other is labelled as Reason R.\\
Assertion A: Global variables increase coupling between modules.\\
Reason R: Modules that share global data are common-coupled, so a change to the data can affect all of them.
\optslong{Both A and R are true and R is the correct explanation of A}{Both A and R are true but R is NOT the correct explanation of A}{A is true but R is false}{A is false but R is true}
\end{question}

\begin{question}{M}
In component-level design, the Common Closure Principle states that
\optslong{all classes should be in one package}{every class needs its own package}{packages must never change}{classes that change together should be packaged together}
\end{question}

\begin{question}{M}
Which of the following is a structural design pattern?
\optsshort{Iterator}{Command}{Observer}{Decorator}
\end{question}

\begin{question}{L}
A good software design should exhibit
\optslong{tight coupling and low cohesion}{loose coupling and high cohesion}{loose coupling and low cohesion}{tight coupling and high cohesion}
\end{question}

\begin{question}{M}
Given below are two statements.\\
Statement I: The pipe-and-filter architecture is well suited to interactive applications with complex user interaction.\\
Statement II: In a layered architecture, each layer may use the services of every other layer, above or below.
\optslong{Both Statement I and Statement II are true}{Both Statement I and Statement II are false}{Statement I is true but Statement II is false}{Statement I is false but Statement II is true}
\end{question}

\begin{answers}
\ans{1}{2}{Coincidental < logical < temporal < procedural < communicational < sequential < functional.}
\ans{2}{2}{Content > common > external > control > stamp > data coupling.}
\ans{3}{3}{Stamp coupling passes a whole data structure when only part is needed.}
\ans{4}{3}{Functional cohesion is the ideal.}
\ans{5}{1}{Cohesion and coupling are the two qualitative measures of functional independence.}
\ans{6}{1}{It limits the propagation of errors and changes.}
\ans{7}{4}{There is a region of minimum cost; too many modules raise integration cost, which is what R says, so A is false.}
\ans{8}{1}{Architectural styles from Pressman and Shaw and Garlan.}
\ans{9}{4}{These, with interface segregation, form the SOLID principles.}
\ans{10}{4}{Patterns from the Gang of Four catalogue.}
\ans{11}{3}{Examples: Builder (creational), Decorator (structural), Strategy (behavioral).}
\ans{12}{3}{It is widely used in web and GUI applications.}
\ans{13}{3}{For example, provide undo, use meaningful defaults and keep conventions consistent.}
\ans{14}{1}{Each step decomposes a high-level statement into more detailed ones.}
\ans{15}{3}{Refactoring improves internal structure without changing external behaviour.}
\ans{16}{1}{Fan-in is the number of modules that call a given module; high fan-in indicates reuse.}
\ans{17}{4}{High fan-in is desirable; very high fan-out suggests that a module controls too many subordinates.}
\ans{18}{4}{Transaction analysis is used for transaction flow.}
\ans{19}{3}{It reduces coupling between clients and subsystem classes.}
\ans{20}{1}{Many client-specific interfaces are better than one general-purpose interface.}
\ans{21}{1}{The client chooses the algorithm at run time.}
\ans{22}{1}{Data, architectural, interface and component-level design are the four design models in Pressman.}
\ans{23}{4}{At the component level, data design chooses data structures.}
\ans{24}{2}{Modularity, aspects and functional independence are manifestations of this concept.}
\ans{25}{1}{Common coupling is the second-worst kind of coupling.}
\ans{26}{4}{Related principles are the Release Reuse Equivalency and Common Reuse principles.}
\ans{27}{4}{Decorator attaches additional responsibilities to an object dynamically.}
\ans{28}{2}{This makes modules independent, understandable and maintainable.}
\ans{29}{2}{Pipe-and-filter suits batch transformation; in a strict layered architecture a layer uses only the layer directly below it.}
\end{answers}

\section{Object-Oriented Analysis and Design and UML}

\begin{question}{H}
Match List I with List II.\\
List I (UML relationship): A. Generalization  B. Composition  C. Aggregation  D. Dependency\\
List II (Notation): I. Dashed line with an open arrowhead  II. Solid line with a hollow triangle pointing to the parent  III. Solid line with a filled diamond at the whole  IV. Solid line with a hollow diamond at the whole
\optsshort{A-II, B-III, C-IV, D-I}{A-II, B-IV, C-III, D-I}{A-III, B-II, C-IV, D-I}{A-II, B-III, C-I, D-IV}
\end{question}

\begin{question}{H}
Given below are two statements.\\
Statement I: In composition, the part cannot exist independently of the whole and is destroyed with it.\\
Statement II: In aggregation, the part may exist independently of the whole.
\optslong{Both Statement I and Statement II are true}{Both Statement I and Statement II are false}{Statement I is true but Statement II is false}{Statement I is false but Statement II is true}
\end{question}

\begin{question}{H}
Match List I with List II.\\
List I (UML diagram): A. Sequence diagram  B. State machine diagram  C. Deployment diagram  D. Use case diagram\\
List II (What it models): I. The physical nodes and the artifacts deployed on them  II. Interactions between objects ordered in time  III. Functional requirements as interactions between actors and the system  IV. The states of an object and the events that cause transitions
\optsshort{A-II, B-I, C-IV, D-III}{A-II, B-IV, C-I, D-III}{A-II, B-IV, C-III, D-I}{A-IV, B-II, C-I, D-III}
\end{question}

\begin{question}{H}
Which of the following are structural (static) UML diagrams?\\
A. Class diagram  B. Component diagram  C. Deployment diagram  D. Activity diagram  E. Object diagram\\
Choose the correct answer from the options given below:
\optsshort{C, D and E only}{A, B, C and E only}{A, B, C, D and E}{A, B and D only}
\end{question}

\begin{question}{H}
In UML class notation, the visibility markers +, $-$ and \# denote respectively
\optslong{protected, private and public}{private, public and protected}{public, protected and private}{public, private and protected}
\end{question}

\begin{question}{M}
UML 2.x defines how many types of diagrams?
\optsshort{12}{16}{14}{9}
\end{question}

\begin{question}{M}
The communication (collaboration) diagram and the sequence diagram are both
\optslong{interaction diagrams showing messages between objects}{structure diagrams}{deployment diagrams}{use case diagrams}
\end{question}

\begin{question}{H}
Given below are two statements: one is labelled as Assertion A and the other is labelled as Reason R.\\
Assertion A: Encapsulation improves maintainability.\\
Reason R: Inheritance lets a subclass reuse the code of its superclass.
\optslong{Both A and R are true and R is the correct explanation of A}{Both A and R are true but R is NOT the correct explanation of A}{A is true but R is false}{A is false but R is true}
\end{question}

\begin{question}{H}
Match List I with List II.\\
List I (CK metric): A. WMC  B. DIT  C. NOC  D. CBO\\
List II (Meaning): I. Number of immediate subclasses of a class  II. Sum of the complexities of the methods of a class  III. Number of other classes to which a class is coupled  IV. Length of the longest path from a class to the root of the inheritance tree
\optsshort{A-II, B-IV, C-I, D-III}{A-II, B-I, C-IV, D-III}{A-IV, B-II, C-I, D-III}{A-II, B-IV, C-III, D-I}
\end{question}

\begin{question}{M}
In the CK metrics suite, LCOM measures
\optslong{lines of code per method}{lack of cohesion among the methods of a class}{the depth of inheritance}{the number of classes}
\end{question}

\begin{question}{M}
The multiplicity 1..* on an association end means
\optsshort{exactly one}{zero or one}{one or more}{zero or more}
\end{question}

\begin{question}{M}
Given below are two statements.\\
Statement I: Inheritance always reduces coupling between a subclass and its superclass.\\
Statement II: Polymorphism allows the same message to produce different behaviour depending on the receiving object's class.
\optslong{Both Statement I and Statement II are true}{Both Statement I and Statement II are false}{Statement I is true but Statement II is false}{Statement I is false but Statement II is true}
\end{question}

\begin{question}{M}
In UML, an ``include'' relationship between use cases is drawn as
\optslong{a solid line}{a dashed arrow labelled \guillemotleft{}include\guillemotright{} from the base use case to the included use case}{a filled diamond}{a hollow triangle}
\end{question}

\begin{question}{M}
The Object Modeling Technique (OMT) was developed by
\optsshort{James Rumbaugh}{Ivar Jacobson}{Peter Coad}{Grady Booch}
\end{question}

\begin{question}{M}
An activity diagram is most similar to
\optslong{a flowchart, showing the flow of activities, decisions and concurrency}{an ER diagram}{a class diagram}{a deployment diagram}
\end{question}

\begin{question}{M}
Given below are two statements: one is labelled as Assertion A and the other is labelled as Reason R.\\
Assertion A: A high depth of inheritance tree (DIT) makes a class harder to understand and test.\\
Reason R: Deeper classes inherit more methods, so their behaviour is harder to predict.
\optslong{Both A and R are true and R is the correct explanation of A}{Both A and R are true but R is NOT the correct explanation of A}{A is true but R is false}{A is false but R is true}
\end{question}

\begin{question}{M}
Which UML diagram shows the organisation of model elements into groups and their dependencies?
\optsshort{Package diagram}{State diagram}{Timing diagram}{Sequence diagram}
\end{question}

\begin{question}{M}
In object-oriented analysis, a class responsible for coordinating a use case's steps is called a
\optsshort{control class}{entity class}{boundary class}{utility class}
\end{question}

\begin{question}{M}
Given below are two statements.\\
Statement I: Abstract classes cannot be instantiated directly.\\
Statement II: An interface in UML is shown with the stereotype \guillemotleft{}interface\guillemotright{} or as a lollipop symbol.
\optslong{Both Statement I and Statement II are true}{Both Statement I and Statement II are false}{Statement I is true but Statement II is false}{Statement I is false but Statement II is true}
\end{question}

\begin{question}{L}
UML is standardised by the
\optslong{W3C}{IEEE}{Object Management Group (OMG)}{ISO only}
\end{question}

\begin{question}{M}
Given below are two statements.\\
Statement I: A sequence diagram shows the static structure of classes and their attributes.\\
Statement II: A state machine diagram is usually drawn for every class, regardless of its behaviour.
\optslong{Both Statement I and Statement II are true}{Both Statement I and Statement II are false}{Statement I is true but Statement II is false}{Statement I is false but Statement II is true}
\end{question}

\begin{answers}
\ans{1}{1}{Realization is a dashed line with a hollow triangle.}
\ans{2}{1}{For example, rooms in a house (composition) versus players in a team (aggregation).}
\ans{3}{2}{UML structural and behavioural diagrams.}
\ans{4}{2}{Activity diagrams are behavioural.}
\ans{5}{4}{\textasciitilde{} denotes package visibility.}
\ans{6}{3}{Seven structure diagrams and seven behaviour diagrams (including four interaction diagrams).}
\ans{7}{1}{The sequence diagram emphasises time ordering; the communication diagram emphasises links between objects.}
\ans{8}{2}{Both are true, but encapsulation helps maintenance by hiding internal data behind an interface, which R does not mention.}
\ans{9}{1}{The Chidamber-Kemerer suite also includes RFC and LCOM.}
\ans{10}{2}{A high LCOM suggests that the class should be split.}
\ans{11}{3}{0..* (or *) means zero or more.}
\ans{12}{4}{Inheritance creates strong coupling between a subclass and its superclass.}
\ans{13}{2}{An \guillemotleft{}extend\guillemotright{} arrow points from the extending use case to the base use case.}
\ans{14}{1}{Booch, Rumbaugh and Jacobson later unified their methods into UML.}
\ans{15}{1}{Swimlanes show which actor or class is responsible for each activity.}
\ans{16}{1}{But a deeper hierarchy also indicates more reuse.}
\ans{17}{1}{Packages group related classes.}
\ans{18}{1}{Boundary classes handle interaction with actors; entity classes hold persistent information.}
\ans{19}{1}{Both are standard UML conventions.}
\ans{20}{3}{The OMG adopted UML in 1997.}
\ans{21}{2}{Sequence diagrams show interactions over time; state diagrams are drawn only for classes with significant state-dependent behaviour.}
\end{answers}

\section{Software Quality, Reliability and Risk Management (McCall, ISO 9126, SQA, CMMI, RMMM)}

\begin{question}{H}
Match List I with List II.\\
List I (McCall's quality factor category): A. Product operation  B. Product revision  C. Product transition\\
List II (Factors): I. Portability, reusability and interoperability  II. Correctness, reliability, efficiency, integrity and usability  III. Maintainability, flexibility and testability
\optsshort{A-II, B-III, C-I}{A-III, B-II, C-I}{A-I, B-III, C-II}{A-II, B-I, C-III}
\end{question}

\begin{question}{H}
Which of the following are the six quality characteristics of ISO 9126?\\
A. Functionality  B. Reliability  C. Usability  D. Efficiency  E. Maintainability  F. Portability\\
Choose the correct answer from the options given below:
\optsshort{A, C, E and F only}{A, B and C only}{A, B, C, D, E and F}{B, D, E and F only}
\end{question}

\begin{question}{H}
Match List I with List II.\\
List I (CMMI maturity level): A. Level 1  B. Level 2  C. Level 3  D. Level 5\\
List II (Name): I. Optimizing  II. Initial  III. Defined  IV. Managed
\optsshort{A-II, B-IV, C-III, D-I}{A-IV, B-II, C-III, D-I}{A-II, B-IV, C-I, D-III}{A-II, B-III, C-IV, D-I}
\end{question}

\begin{question}{H}
A system has a mean time to failure (MTTF) of 68 days and a mean time to repair (MTTR) of 3 days. Its availability is approximately
\optsshort{97.1\%}{90.4\%}{99.0\%}{95.8\%}
\end{question}

\begin{question}{H}
Given below are two statements.\\
Statement I: MTBF = MTTF + MTTR.\\
Statement II: Availability is independent of the mean time to repair.
\optslong{Both Statement I and Statement II are true}{Both Statement I and Statement II are false}{Statement I is true but Statement II is false}{Statement I is false but Statement II is true}
\end{question}

\begin{question}{H}
A team finds 90 errors before delivery, and users report 10 defects after delivery. The defect removal efficiency (DRE) is
\optsshort{0.90}{0.81}{0.10}{0.99}
\end{question}

\begin{question}{H}
Given below are two statements: one is labelled as Assertion A and the other is labelled as Reason R.\\
Assertion A: Quality assurance focuses on the process, while quality control focuses on the product.\\
Reason R: Quality assurance aims to prevent defects, while quality control aims to detect them.
\optslong{Both A and R are true and R is the correct explanation of A}{Both A and R are true but R is NOT the correct explanation of A}{A is true but R is false}{A is false but R is true}
\end{question}

\begin{question}{H}
A risk has a probability of 70\% and an estimated cost (impact) of Rs. 50,000 if it occurs. The risk exposure is
\optsshort{Rs. 35,000}{Rs. 50,000}{Rs. 70,000}{Rs. 15,000}
\end{question}

\begin{question}{H}
Match List I with List II.\\
List I (Risk category): A. Project risks  B. Technical risks  C. Business risks  D. Unpredictable risks\\
List II (Example): I. Building a product that no one wants (market risk)  II. Budget overrun and staff turnover  III. Risks that are extremely difficult to identify in advance  IV. Uncertainty about the interface or implementation technology
\optsshort{A-II, B-IV, C-I, D-III}{A-II, B-IV, C-III, D-I}{A-IV, B-II, C-I, D-III}{A-II, B-I, C-IV, D-III}
\end{question}

\begin{question}{H}
RMMM stands for
\optslong{Reliability Measurement and Monitoring Method}{Risk Mitigation, Monitoring and Management}{Requirement Modelling and Maintenance Management}{Risk Measurement, Modelling and Mapping}
\end{question}

\begin{question}{M}
Given below are two statements.\\
Statement I: Reactive risk strategies deal with risks only when problems occur.\\
Statement II: Proactive risk strategies identify, assess and plan for risks before technical work begins.
\optslong{Both Statement I and Statement II are true}{Both Statement I and Statement II are false}{Statement I is true but Statement II is false}{Statement I is false but Statement II is true}
\end{question}

\begin{question}{M}
A risk table (risk projection) lists for each risk its
\optslong{test cases}{category, probability and impact}{design patterns}{cost per line of code}
\end{question}

\begin{question}{M}
Match List I with List II.\\
List I (Cost of quality component): A. Prevention costs  B. Appraisal costs  C. Internal failure costs  D. External failure costs\\
List II (Example): I. Rework due to defects found before delivery  II. Training, planning and quality planning  III. Complaint resolution and warranty work after delivery  IV. Inspections, testing and equipment calibration
\optsshort{A-IV, B-II, C-I, D-III}{A-II, B-I, C-IV, D-III}{A-II, B-IV, C-III, D-I}{A-II, B-IV, C-I, D-III}
\end{question}

\begin{question}{M}
Six Sigma aims at a defect rate of no more than
\optslong{3.4 defects per million opportunities}{34 defects per million}{6 defects per hundred}{1 defect per thousand}
\end{question}

\begin{question}{M}
A formal technical review (FTR) typically involves
\optslong{3 to 5 people, advance preparation of no more than about 2 hours, and a meeting of less than 2 hours}{only the author}{no records}{20 to 30 people for a full day}
\end{question}

\begin{question}{M}
Given below are two statements: one is labelled as Assertion A and the other is labelled as Reason R.\\
Assertion A: Software reliability can be improved by testing alone, without good design.\\
Reason R: Testing can show the presence of defects but cannot prove their absence.
\optslong{Both A and R are true and R is the correct explanation of A}{Both A and R are true but R is NOT the correct explanation of A}{A is true but R is false}{A is false but R is true}
\end{question}

\begin{question}{M}
Fagan inspection is
\optslong{a risk-analysis method}{a formal, highly structured review process with defined roles such as moderator, reader and recorder}{a type of unit testing}{an informal walkthrough}
\end{question}

\begin{question}{M}
The ISO 9001 standard applies to
\optslong{programming languages}{quality management systems of organisations in any industry, including software}{only software products}{network protocols}
\end{question}

\begin{question}{M}
Given below are two statements.\\
Statement I: Statistical SQA ignores the causes of defects.\\
Statement II: Statistical SQA uses the Pareto principle: about 80\% of defects come from about 20\% of causes.
\optslong{Both Statement I and Statement II are true}{Both Statement I and Statement II are false}{Statement I is true but Statement II is false}{Statement I is false but Statement II is true}
\end{question}

\begin{question}{M}
The probability of failure on demand (POFOD) is an appropriate reliability metric for
\optslong{batch payroll systems}{protection systems that are called upon only occasionally, such as an emergency shutdown}{games}{continuously running systems}
\end{question}

\begin{question}{M}
In risk management, the risk referent level is
\optslong{the probability of a risk}{the point at which the project should be stopped or significantly changed if a risk (or combination of risks) reaches it}{the number of risks}{the cost of mitigation}
\end{question}

\begin{question}{M}
Given below are two statements.\\
Statement I: CMMI level 3 (Defined) requires that the organisation's processes be documented, standardised and integrated.\\
Statement II: At CMMI level 1 (Initial), processes are well defined and predictable.
\optslong{Both Statement I and Statement II are true}{Both Statement I and Statement II are false}{Statement I is true but Statement II is false}{Statement I is false but Statement II is true}
\end{question}

\begin{question}{L}
Software reliability is defined as
\optslong{the number of lines of code}{the probability of failure-free operation of a program in a specified environment for a specified time}{the cost of development}{the speed of execution}
\end{question}

\begin{question}{M}
Musa's basic execution time model and the Jelinski-Moranda model are examples of
\optslong{design patterns}{software reliability growth models}{process models}{cost estimation models}
\end{question}

\begin{answers}
\ans{1}{1}{McCall's model has 11 quality factors in three groups.}
\ans{2}{3}{ISO/IEC 25010 replaced ISO 9126 and adds security and compatibility.}
\ans{3}{1}{Level 4 is Quantitatively Managed (called Managed in the original CMM, where level 2 was Repeatable).}
\ans{4}{4}{Availability = MTTF / (MTTF + MTTR) = 68 / 71 $\approx$ 0.958.}
\ans{5}{3}{Availability = MTTF / (MTTF + MTTR), so it falls as repair time grows.}
\ans{6}{1}{DRE = E / (E + D) = 90 / 100.}
\ans{7}{1}{QA defines and audits processes; QC inspects and tests work products.}
\ans{8}{1}{RE = P $\times$ C = 0.7 $\times$ 50,000.}
\ans{9}{1}{Pressman classifies risks as project, technical and business, and as known, predictable and unpredictable.}
\ans{10}{2}{The RMMM plan documents all risk-analysis work.}
\ans{11}{1}{Pressman describes reactive strategies as ``fire-fighting mode''.}
\ans{12}{2}{Risks are sorted by probability and impact, and a cut-off line separates those that will be managed.}
\ans{13}{4}{External failure costs are the most expensive.}
\ans{14}{1}{Its improvement methodology is DMAIC: define, measure, analyse, improve and control.}
\ans{15}{1}{The review focuses on the product, not the producer.}
\ans{16}{4}{R (Dijkstra) is true; it is exactly why reliability must also be built in through good processes and design, so A is false.}
\ans{17}{2}{Inspections are among the most effective defect-detection techniques.}
\ans{18}{2}{ISO 90003 gives guidelines for applying ISO 9001 to software.}
\ans{19}{4}{Statistical SQA collects and categorises defect data to find and fix the vital few causes.}
\ans{20}{2}{ROCOF (rate of occurrence of failures) suits continuously running systems.}
\ans{21}{2}{Typical referents are cost, schedule and performance.}
\ans{22}{3}{At level 1, processes are ad hoc and success depends on individual effort.}
\ans{23}{2}{It is one of the most important quality attributes.}
\ans{24}{2}{They predict how reliability improves as defects are removed during testing.}
\end{answers}

\section{Estimation and Scheduling (LOC and FP estimation, COCOMO, Putnam, scheduling, earned value)}

\begin{question}{H}
Using the basic COCOMO model, the effort for an organic-mode project of 32 KLOC (E = 2.4 $\times$ KLOC\textsuperscript{1.05} person-months) is approximately
\optsshort{91.3 PM}{230.4 PM}{145.5 PM}{76.8 PM}
\end{question}

\begin{question}{H}
For an organic-mode project with an effort of 91.3 person-months, the basic COCOMO development time D = 2.5 $\times$ E\textsuperscript{0.38} is approximately
\optsshort{36.5 months}{13.9 months}{9.1 months}{18.2 months}
\end{question}

\begin{question}{H}
Match List I with List II.\\
List I (Basic COCOMO mode): A. Organic  B. Semi-detached  C. Embedded\\
List II (Effort equation): I. E = 3.6 (KLOC)\textsuperscript{1.20}  II. E = 2.4 (KLOC)\textsuperscript{1.05}  III. E = 3.0 (KLOC)\textsuperscript{1.12}
\optsshort{A-III, B-II, C-I}{A-II, B-I, C-III}{A-I, B-III, C-II}{A-II, B-III, C-I}
\end{question}

\begin{question}{H}
Given below are two statements.\\
Statement I: An embedded-mode COCOMO project operates under tight hardware, software and operational constraints.\\
Statement II: For the same size, basic COCOMO gives the least effort for embedded-mode projects.
\optslong{Both Statement I and Statement II are true}{Both Statement I and Statement II are false}{Statement I is true but Statement II is false}{Statement I is false but Statement II is true}
\end{question}

\begin{question}{H}
A system has 20 external inputs, 15 external outputs, 10 external inquiries, 5 internal logical files and 3 external interface files, all of average complexity (weights 4, 5, 4, 10 and 7). The unadjusted function point count is
\optsshort{300}{342}{266}{228}
\end{question}

\begin{question}{H}
A system has 266 unadjusted function points, and the sum of the 14 value adjustment factors ($\Sigma$Fi) is 42. The adjusted function point count is approximately
\optsshort{308}{359}{285}{266}
\end{question}

\begin{question}{H}
The value adjustment factor in function point analysis, 0.65 + 0.01 $\times$ $\Sigma$Fi with 14 factors each rated 0 to 5, lies in the range
\optsshort{0.5 to 1.5}{0.65 to 1.35}{1 to 2}{0 to 1}
\end{question}

\begin{question}{H}
Using three-point estimation, a module has optimistic, most likely and pessimistic sizes of 4600, 6900 and 8600 LOC. The expected size is
\optsshort{6700 LOC}{6900 LOC}{7000 LOC}{6800 LOC}
\end{question}

\begin{question}{H}
Match List I with List II.\\
List I (Function point component): A. External input  B. External output  C. External inquiry  D. Internal logical file\\
List II (Description): I. A logical group of data maintained within the application boundary  II. Data or control information entering the application from outside  III. An input-output combination that retrieves data without derived data  IV. Derived information, such as reports, sent outside the application
\optsshort{A-II, B-III, C-IV, D-I}{A-IV, B-II, C-III, D-I}{A-II, B-IV, C-I, D-III}{A-II, B-IV, C-III, D-I}
\end{question}

\begin{question}{H}
Given below are two statements: one is labelled as Assertion A and the other is labelled as Reason R.\\
Assertion A: Function points can be estimated early in the life cycle, before the code is written.\\
Reason R: Function points are derived from the requirements (inputs, outputs, inquiries and files), and are independent of the programming language.
\optslong{Both A and R are true and R is the correct explanation of A}{Both A and R are true but R is NOT the correct explanation of A}{A is true but R is false}{A is false but R is true}
\end{question}

\begin{question}{H}
Brooks' law states that
\optslong{adding manpower to a late software project makes it later}{work expands to fill the time available}{software projects are always on time}{the cost of fixing a defect doubles every phase}
\end{question}

\begin{question}{M}
The number of communication paths among a team of 10 developers is
\optsshort{45}{10}{100}{20}
\end{question}

\begin{question}{H}
In earned value analysis, BCWS = 100, BCWP = 80 and ACWP = 90 (in thousands). The schedule performance index (SPI) and cost performance index (CPI) are
\optsshort{1.25 and 1.13}{0.8 and 1.13}{0.8 and 0.89}{0.89 and 0.8}
\end{question}

\begin{question}{M}
In earned value analysis, the schedule variance (SV) and cost variance (CV) are defined as
\optsshort{BCWS $-$ BCWP and ACWP $-$ BCWP}{ACWP $-$ BCWS and BCWP $-$ ACWP}{BCWS $-$ ACWP and BCWP $-$ BCWS}{BCWP $-$ BCWS and BCWP $-$ ACWP}
\end{question}

\begin{question}{H}
In a project network, activity A (3 days) precedes B (4 days) and C (2 days); B precedes D (5 days); C precedes E (6 days); D and E precede F (2 days). The length of the critical path is
\optsshort{16 days}{12 days}{14 days}{13 days}
\end{question}

\begin{question}{M}
In the project network where A (3) precedes B (4) and C (2), B precedes D (5), C precedes E (6), and D and E precede F (2), the slack (float) of activity C is
\optsshort{0 days}{3 days}{2 days}{1 day}
\end{question}

\begin{question}{M}
The 40-20-40 rule of effort distribution assigns
\optslong{40\% to coding}{40\% to maintenance}{40\% to front-end analysis and design, 20\% to coding and 40\% to testing}{equal effort to all phases}
\end{question}

\begin{question}{M}
Given below are two statements.\\
Statement I: In the Putnam (SLIM) model, effort is very sensitive to the delivery time, varying inversely with its fourth power.\\
Statement II: The Putnam model is based on the Rayleigh-Norden curve of staffing.
\optslong{Both Statement I and Statement II are true}{Both Statement I and Statement II are false}{Statement I is true but Statement II is false}{Statement I is false but Statement II is true}
\end{question}

\begin{question}{M}
Intermediate COCOMO refines basic COCOMO by multiplying the nominal effort by
\optslong{the number of function points}{the number of developers}{a constant of 2}{an effort adjustment factor computed from 15 cost drivers}
\end{question}

\begin{question}{M}
COCOMO II defines three sub-models, namely
\optslong{waterfall, spiral and agile}{organic, semi-detached and embedded}{basic, intermediate and detailed}{application composition, early design and post-architecture}
\end{question}

\begin{question}{M}
A project of 33,200 LOC is estimated with an average productivity of 620 LOC per person-month. The estimated effort is about
\optsshort{80 PM}{62 PM}{54 PM}{33 PM}
\end{question}

\begin{question}{M}
Given below are two statements: one is labelled as Assertion A and the other is labelled as Reason R.\\
Assertion A: Lines-of-code measures are language-dependent.\\
Reason R: Function points are measured in lines of code.
\optslong{Both A and R are true and R is the correct explanation of A}{Both A and R are true but R is NOT the correct explanation of A}{A is true but R is false}{A is false but R is true}
\end{question}

\begin{question}{M}
A Gantt chart (timeline chart) shows
\optslong{the tasks of a project as horizontal bars against a time axis}{the data flow}{the class structure}{the risk table}
\end{question}

\begin{question}{M}
PERT estimates the expected duration of an activity as
\optsshort{(o + m + p)/3}{(o + 4m + p)/6}{(o + 2m + p)/4}{m}
\end{question}

\begin{question}{M}
The wideband Delphi technique is used for
\optslong{expert-judgment-based estimation through anonymous rounds of estimates and discussion}{configuration management}{coding}{testing}
\end{question}

\begin{question}{M}
For a vocabulary of n1 = 10 distinct operators and n2 = 7 distinct operands, with N1 = 16 and N2 = 15 total occurrences, the Halstead program volume V = N log2 n is approximately
\optsshort{155}{180}{127}{100}
\end{question}

\begin{question}{M}
Given below are two statements.\\
Statement I: Halstead's program length N is the total number of operator and operand occurrences.\\
Statement II: Halstead's vocabulary n is the number of distinct operators plus the number of distinct operands.
\optslong{Both Statement I and Statement II are true}{Both Statement I and Statement II are false}{Statement I is true but Statement II is false}{Statement I is false but Statement II is true}
\end{question}

\begin{question}{M}
Given below are two statements.\\
Statement I: Activities on the critical path have positive slack.\\
Statement II: A work breakdown structure (WBS) divides the project into a hierarchy of smaller, manageable tasks.
\optslong{Both Statement I and Statement II are true}{Both Statement I and Statement II are false}{Statement I is true but Statement II is false}{Statement I is false but Statement II is true}
\end{question}

\begin{question}{L}
KLOC stands for
\optslong{key lines of code}{known lines of code}{kernel lines of code}{kilo (thousand) lines of code}
\end{question}

\begin{question}{M}
Given below are two statements.\\
Statement I: In basic COCOMO, effort grows linearly with size for every mode.\\
Statement II: In basic COCOMO, development time grows faster than effort.
\optslong{Both Statement I and Statement II are true}{Both Statement I and Statement II are false}{Statement I is true but Statement II is false}{Statement I is false but Statement II is true}
\end{question}

\begin{answers}
\ans{1}{1}{2.4 $\times$ 32\textsuperscript{1.05} $\approx$ 91.3 person-months.}
\ans{2}{2}{2.5 $\times$ 91.3\textsuperscript{0.38} $\approx$ 13.9 months, so about 91.3 / 13.9 $\approx$ 6.6 people on average.}
\ans{3}{4}{For a 32 KLOC project these give about 91, 146 and 230 person-months.}
\ans{4}{3}{Embedded mode has the largest coefficient and exponent, and hence the most effort.}
\ans{5}{3}{20 $\times$ 4 + 15 $\times$ 5 + 10 $\times$ 4 + 5 $\times$ 10 + 3 $\times$ 7 = 80 + 75 + 40 + 50 + 21 = 266.}
\ans{6}{3}{FP = UFP $\times$ (0.65 + 0.01 $\times$ $\Sigma$Fi) = 266 $\times$ 1.07 $\approx$ 284.6.}
\ans{7}{2}{$\Sigma$Fi ranges from 0 to 70.}
\ans{8}{4}{(4600 + 4 $\times$ 6900 + 8600) / 6 = 6800 LOC.}
\ans{9}{4}{External interface files are maintained by other applications.}
\ans{10}{1}{LOC can be known accurately only after coding.}
\ans{11}{1}{New people need training, and communication overhead grows.}
\ans{12}{1}{n(n $-$ 1)/2 = 45; this growth in paths is one reason behind Brooks' law.}
\ans{13}{3}{SPI = BCWP / BCWS = 0.8 (behind schedule); CPI = BCWP / ACWP $\approx$ 0.89 (over budget).}
\ans{14}{4}{Negative values indicate being behind schedule or over budget.}
\ans{15}{3}{A--B--D--F takes 3 + 4 + 5 + 2 = 14 days; A--C--E--F takes 13 days.}
\ans{16}{4}{C lies on the path A--C--E--F of 13 days, one day shorter than the critical path.}
\ans{17}{3}{It is a rough guideline for distributing development effort.}
\ans{18}{1}{Compressing the schedule therefore increases effort sharply.}
\ans{19}{4}{Cost drivers cover product, hardware, personnel and project attributes.}
\ans{20}{4}{The application composition model uses object points.}
\ans{21}{3}{33,200 / 620 $\approx$ 53.5 person-months.}
\ans{22}{3}{Function points are language-independent measures of functionality; LOC depends on the language because the same function takes different numbers of lines.}
\ans{23}{1}{PERT/CPM networks show the dependencies between tasks.}
\ans{24}{2}{o is optimistic, m most likely and p pessimistic time.}
\ans{25}{1}{Estimates converge over several rounds.}
\ans{26}{3}{N = 31 and n = 17, so V = 31 $\times$ log2 17 $\approx$ 126.7.}
\ans{27}{1}{Both are basic Halstead measures.}
\ans{28}{4}{Critical-path activities have zero slack; any delay in them delays the project.}
\ans{29}{4}{Effort is often expressed in person-months per KLOC.}
\ans{30}{2}{The exponents above 1 make effort grow faster than size, and time grows much more slowly than effort.}
\end{answers}

\section{Software Testing (V\&V, white-box and black-box testing, basis path testing, test levels)}

\begin{question}{H}
A program's control flow graph has 9 edges and 7 nodes. Its cyclomatic complexity is
\optsshort{4}{3}{2}{5}
\end{question}

\begin{question}{H}
The code if (a > b) \{ if (b > c) x = 1; else x = 2; \} else x = 3; while (i < n) i++; has cyclomatic complexity
\optsshort{2}{5}{3}{4}
\end{question}

\begin{question}{H}
Which of the following are valid ways to compute cyclomatic complexity V(G) of a flow graph?\\
A. E $-$ N + 2  B. Number of predicate nodes + 1  C. Number of regions of the planar flow graph  D. Number of nodes $-$ 1\\
Choose the correct answer from the options given below:
\optsshort{A, B and C only}{B, C and D only}{A, B, C and D}{A and D only}
\end{question}

\begin{question}{H}
Given below are two statements.\\
Statement I: Cyclomatic complexity gives an upper bound on the number of tests needed to execute every independent path at least once.\\
Statement II: Cyclomatic complexity gives the total number of possible paths through the program.
\optslong{Both Statement I and Statement II are true}{Both Statement I and Statement II are false}{Statement I is true but Statement II is false}{Statement I is false but Statement II is true}
\end{question}

\begin{question}{H}
Match List I with List II.\\
List I (Term): A. Error  B. Fault (defect, bug)  C. Failure  D. Test oracle\\
List II (Meaning): I. A mechanism to decide whether a test passed  II. A human mistake  III. The manifestation of a fault as incorrect behaviour during execution  IV. An incorrect step or data definition in the program
\optsshort{A-IV, B-II, C-III, D-I}{A-II, B-III, C-IV, D-I}{A-II, B-IV, C-I, D-III}{A-II, B-IV, C-III, D-I}
\end{question}

\begin{question}{H}
Given below are two statements.\\
Statement I: Verification asks ``Are we building the product right?''\\
Statement II: Validation asks ``Are we building the right product?''
\optslong{Both Statement I and Statement II are true}{Both Statement I and Statement II are false}{Statement I is true but Statement II is false}{Statement I is false but Statement II is true}
\end{question}

\begin{question}{H}
Match List I with List II.\\
List I (Testing technique): A. Equivalence partitioning  B. Boundary value analysis  C. Basis path testing  D. Cause-effect graphing\\
List II (Type and idea): I. White-box; uses cyclomatic complexity to derive independent paths  II. Black-box; tests values at the edges of input domains  III. Black-box; maps combinations of input conditions to actions via a decision table  IV. Black-box; divides inputs into classes that are treated alike
\optsshort{A-IV, B-II, C-III, D-I}{A-IV, B-II, C-I, D-III}{A-IV, B-I, C-II, D-III}{A-II, B-IV, C-I, D-III}
\end{question}

\begin{question}{H}
For a function with 2 independent input variables, the number of test cases needed by basic boundary value analysis (min, min+, nominal, max$-$, max for each variable, holding the others at nominal) is
\optsshort{9}{13}{5}{25}
\end{question}

\begin{question}{M}
For a function with 3 input variables, worst-case boundary value testing needs how many test cases?
\optsshort{125}{13}{343}{19}
\end{question}

\begin{question}{H}
Arrange the levels of testing in the order in which they are normally performed.\\
A. System testing  B. Unit testing  C. Acceptance testing  D. Integration testing
\optsshort{B, D, A, C}{B, A, D, C}{D, B, A, C}{B, D, C, A}
\end{question}

\begin{question}{H}
Match List I with List II.\\
List I (Test type): A. Alpha testing  B. Beta testing  C. Regression testing  D. Stress testing\\
List II (Description): I. Re-running tests after changes to ensure nothing else broke  II. Conducted at the developer's site by end users  III. Executing the system under abnormal quantities, frequencies or volumes of resources  IV. Conducted at one or more customer sites, without the developer present
\optsshort{A-IV, B-II, C-I, D-III}{A-II, B-I, C-IV, D-III}{A-II, B-IV, C-III, D-I}{A-II, B-IV, C-I, D-III}
\end{question}

\begin{question}{H}
Given below are two statements: one is labelled as Assertion A and the other is labelled as Reason R.\\
Assertion A: In top-down integration testing, drivers are needed to call the top-level module.\\
Reason R: Lower-level modules called by the module under test may not yet be ready.
\optslong{Both A and R are true and R is the correct explanation of A}{Both A and R are true but R is NOT the correct explanation of A}{A is true but R is false}{A is false but R is true}
\end{question}

\begin{question}{M}
In bottom-up integration testing, the components that simulate higher-level modules calling the module under test are called
\optsshort{oracles}{drivers}{stubs}{mocks only}
\end{question}

\begin{question}{H}
Given below are two statements.\\
Statement I: Statement coverage subsumes branch coverage.\\
Statement II: Branch (decision) coverage subsumes statement coverage.
\optslong{Both Statement I and Statement II are true}{Both Statement I and Statement II are false}{Statement I is true but Statement II is false}{Statement I is false but Statement II is true}
\end{question}

\begin{question}{M}
Arrange the following coverage criteria from the weakest to the strongest.\\
A. Path coverage  B. Statement coverage  C. Branch coverage
\optsshort{B, A, C}{C, B, A}{A, C, B}{B, C, A}
\end{question}

\begin{question}{M}
The statement ``Program testing can be used to show the presence of bugs, but never to show their absence'' is due to
\optsshort{Edsger Dijkstra}{Watts Humphrey}{Barry Boehm}{Glenford Myers}
\end{question}

\begin{question}{M}
In mutation testing, 50 mutants are generated, 5 of them are equivalent to the original program, and the test suite kills 40. The mutation score is about
\optsshort{85\%}{80\%}{89\%}{90\%}
\end{question}

\begin{question}{M}
Smoke testing is
\optslong{testing done by customers}{a preliminary integration test, often daily, to check that the build's critical functions work}{a test of the system under extreme heat}{testing only the user interface}
\end{question}

\begin{question}{M}
Data-flow testing selects test paths according to
\optslong{the user interface}{the number of loops}{the locations of definitions and uses of variables (DU chains)}{equivalence classes}
\end{question}

\begin{question}{M}
Given below are two statements: one is labelled as Assertion A and the other is labelled as Reason R.\\
Assertion A: Black-box testing can detect missing functionality.\\
Reason R: Black-box tests are derived from the code structure.
\optslong{Both A and R are true and R is the correct explanation of A}{Both A and R are true but R is NOT the correct explanation of A}{A is true but R is false}{A is false but R is true}
\end{question}

\begin{question}{M}
Loop testing distinguishes among
\optslong{infinite and finite loops}{recursive and iterative loops}{simple, nested, concatenated and unstructured loops}{for, while and do-while loops only}
\end{question}

\begin{question}{M}
Recovery testing checks
\optslong{how well the system recovers from crashes, hardware failures and similar problems}{the speed of the system}{the user interface}{the documentation}
\end{question}

\begin{question}{M}
Given below are two statements.\\
Statement I: Static testing, such as reviews and inspections, finds defects without executing the program.\\
Statement II: Dynamic testing requires executing the program.
\optslong{Both Statement I and Statement II are true}{Both Statement I and Statement II are false}{Statement I is true but Statement II is false}{Statement I is false but Statement II is true}
\end{question}

\begin{question}{M}
A flow graph has 12 edges, 10 nodes and 2 connected components. Its cyclomatic complexity is
\optsshort{8}{2}{4}{6}
\end{question}

\begin{question}{M}
The debugging approach that works backwards from the place where a symptom appears to the cause is called
\optsshort{backtracking}{cause elimination}{regression}{brute force}
\end{question}

\begin{question}{M}
Given below are two statements.\\
Statement I: Performance testing measures run-time performance in the context of an integrated system.\\
Statement II: Security testing verifies that protection mechanisms built into a system protect it from improper penetration.
\optslong{Both Statement I and Statement II are true}{Both Statement I and Statement II are false}{Statement I is true but Statement II is false}{Statement I is false but Statement II is true}
\end{question}

\begin{question}{L}
Unit testing is usually performed by
\optsshort{the end users}{the developers of the module}{the customers at their site}{the marketing team}
\end{question}

\begin{question}{M}
Given below are two statements.\\
Statement I: Exhaustive testing of all inputs is practical for most programs.\\
Statement II: White-box testing alone can guarantee that a program meets its specification.
\optslong{Both Statement I and Statement II are true}{Both Statement I and Statement II are false}{Statement I is true but Statement II is false}{Statement I is false but Statement II is true}
\end{question}

\begin{question}{M}
A decision table with 3 independent Boolean conditions has at most how many rules (columns)?
\optsshort{6}{9}{8}{3}
\end{question}

\begin{answers}
\ans{1}{1}{V(G) = E $-$ N + 2 = 9 $-$ 7 + 2 = 4.}
\ans{2}{4}{There are three decisions (two ifs and one while): V(G) = 3 + 1 = 4.}
\ans{3}{1}{All three give the same value for a connected flow graph.}
\ans{4}{3}{It counts linearly independent (basis) paths; the number of complete paths can be far larger, or infinite with loops.}
\ans{5}{4}{An error leads to a fault, which may cause a failure.}
\ans{6}{1}{Boehm's well-known distinction.}
\ans{7}{2}{Black-box techniques use the specification; white-box techniques use the code.}
\ans{8}{1}{4n + 1 = 9; robustness testing needs 6n + 1 = 13 and worst-case testing 5\textsuperscript{n} = 25.}
\ans{9}{1}{5\textsuperscript{n} = 5\textsuperscript{3} = 125 (robust worst case needs 7\textsuperscript{3} = 343).}
\ans{10}{1}{Each level tests larger units of the system.}
\ans{11}{4}{Types of testing from Pressman.}
\ans{12}{4}{Top-down integration needs stubs for missing lower-level modules (as R implies); drivers are needed in bottom-up integration.}
\ans{13}{2}{Stubs replace lower-level modules called by the unit under test.}
\ans{14}{4}{Covering every branch executes every statement, but not the other way round (for example an if with no else).}
\ans{15}{4}{Full path coverage is usually infeasible with loops.}
\ans{16}{1}{Exhaustive testing is impossible for any non-trivial program.}
\ans{17}{3}{Mutation score = killed / (total $-$ equivalent) = 40 / 45 $\approx$ 0.89.}
\ans{18}{2}{It is common in projects using daily builds.}
\ans{19}{3}{It is a white-box control-structure testing technique.}
\ans{20}{3}{Black-box tests come from the specification, which is exactly why they can reveal functions the code omits.}
\ans{21}{3}{Unstructured loops should be redesigned before testing.}
\ans{22}{1}{It forces the software to fail and verifies that recovery is performed properly.}
\ans{23}{1}{Both are forms of verification.}
\ans{24}{4}{V(G) = E $-$ N + 2P = 12 $-$ 10 + 4 = 6.}
\ans{25}{1}{Cause elimination uses induction or deduction with binary partitioning.}
\ans{26}{1}{Both are kinds of system testing.}
\ans{27}{2}{It focuses on the smallest unit of design, such as a component or class.}
\ans{28}{2}{Exhaustive testing is infeasible, and white-box testing cannot detect missing functions required by the specification.}
\ans{29}{3}{Each condition can be true or false: 2\textsuperscript{3} = 8 combinations.}
\end{answers}

\section{Software Configuration Management, Maintenance, Reuse, Re-engineering and Reverse Engineering}

\begin{question}{H}
Match List I with List II.\\
List I (Maintenance type): A. Corrective  B. Adaptive  C. Perfective  D. Preventive\\
List II (Purpose): I. Changing the software to improve future maintainability or reliability  II. Fixing defects discovered after delivery  III. Adding enhancements and improving performance as user requirements evolve  IV. Modifying the software to work in a changed environment, such as a new OS
\optsshort{A-II, B-IV, C-I, D-III}{A-IV, B-II, C-III, D-I}{A-II, B-IV, C-III, D-I}{A-II, B-III, C-IV, D-I}
\end{question}

\begin{question}{H}
Which type of maintenance usually accounts for the largest share of maintenance effort?
\optsshort{Corrective}{Perfective}{Adaptive}{Preventive}
\end{question}

\begin{question}{H}
Given below are two statements.\\
Statement I: A baseline is a specification or product that has been formally reviewed and agreed upon and can be changed only through formal change control.\\
Statement II: A software configuration item (SCI) is information created as part of the software process, such as a document, test suite or program component.
\optslong{Both Statement I and Statement II are true}{Both Statement I and Statement II are false}{Statement I is true but Statement II is false}{Statement I is false but Statement II is true}
\end{question}

\begin{question}{H}
Which of the following are tasks in software configuration management?\\
A. Identification of configuration items  B. Version control  C. Change control  D. Configuration auditing  E. Status reporting\\
Choose the correct answer from the options given below:
\optsshort{A, B, C, D and E}{A, B and C only}{B, C and D only}{A, D and E only}
\end{question}

\begin{question}{H}
In change control, the body that evaluates and approves or rejects change requests is the
\optslong{quality circle}{test team}{change control board (or change control authority)}{end users only}
\end{question}

\begin{question}{H}
Arrange the activities of Pressman's software re-engineering process model in a typical order.\\
A. Reverse engineering  B. Inventory analysis  C. Forward engineering  D. Document restructuring  E. Code restructuring
\optsshort{A, B, D, E, C}{B, D, A, E, C}{D, B, A, C, E}{B, A, D, E, C}
\end{question}

\begin{question}{H}
Given below are two statements: one is labelled as Assertion A and the other is labelled as Reason R.\\
Assertion A: Reverse engineering extracts design and specification information from existing source code.\\
Reason R: Reverse engineering changes the functionality of the system.
\optslong{Both A and R are true and R is the correct explanation of A}{Both A and R are true but R is NOT the correct explanation of A}{A is true but R is false}{A is false but R is true}
\end{question}

\begin{question}{M}
Forward engineering in the context of re-engineering means
\optslong{testing only}{analysing old code}{deleting the old system}{using recovered design information to restructure or rebuild the system with modern methods}
\end{question}

\begin{question}{M}
Given below are two statements.\\
Statement I: Git is a distributed version control system in which every developer has a full copy of the repository history.\\
Statement II: Subversion (SVN) is a distributed version control system.
\optslong{Both Statement I and Statement II are true}{Both Statement I and Statement II are false}{Statement I is true but Statement II is false}{Statement I is false but Statement II is true}
\end{question}

\begin{question}{M}
Match List I with List II.\\
List I (Lehman's law of software evolution): A. Continuing change  B. Increasing complexity  C. Continuing growth  D. Declining quality\\
List II (Statement): I. Functionality must grow to maintain user satisfaction  II. A system must be continually adapted or it becomes progressively less useful  III. Quality appears to decline unless the system is adapted to environmental changes  IV. As a system evolves, its complexity increases unless work is done to reduce it
\optsshort{A-II, B-IV, C-III, D-I}{A-II, B-I, C-IV, D-III}{A-II, B-IV, C-I, D-III}{A-IV, B-II, C-I, D-III}
\end{question}

\begin{question}{M}
Maintenance typically accounts for what share of the total life-cycle cost of software?
\optsshort{About 25\%}{About 60\% or more}{Less than 5\%}{About 10\%}
\end{question}

\begin{question}{M}
The abstraction level of reverse engineering refers to
\optslong{the testing level}{the number of lines of code}{the sophistication of the design information that can be extracted, from procedural design up to data flow and ER models}{the programming language}
\end{question}

\begin{question}{M}
Given below are two statements: one is labelled as Assertion A and the other is labelled as Reason R.\\
Assertion A: Software reuse can reduce development cost and improve reliability.\\
Reason R: Reused components have already been tested and used in working systems.
\optslong{Both A and R are true and R is the correct explanation of A}{Both A and R are true but R is NOT the correct explanation of A}{A is true but R is false}{A is false but R is true}
\end{question}

\begin{question}{M}
COTS stands for
\optslong{Common Object Transfer Standard}{Commercial Off-The-Shelf}{Code Optimization Tool Set}{Configuration Of Test Suites}
\end{question}

\begin{question}{M}
Domain engineering in software reuse involves
\optslong{fixing defects}{identifying, constructing, cataloguing and disseminating reusable components for an application domain}{writing user manuals}{testing one program}
\end{question}

\begin{question}{M}
Given below are two statements.\\
Statement I: Legacy systems are always easy to maintain because they are well documented.\\
Statement II: A legacy system is an old system that remains in use because it is critical to the business.
\optslong{Both Statement I and Statement II are true}{Both Statement I and Statement II are false}{Statement I is true but Statement II is false}{Statement I is false but Statement II is true}
\end{question}

\begin{question}{M}
Code restructuring in re-engineering
\optslong{modifies the code to improve its structure and quality without changing its function}{translates the code to another hardware platform only}{removes all comments}{changes the external behaviour of the code}
\end{question}

\begin{question}{M}
In version control, branching is used to
\optslong{compile the code}{delete old versions}{develop features or fixes in parallel without affecting the main line, merging them later}{test the system}
\end{question}

\begin{question}{M}
Given below are two statements: one is labelled as Assertion A and the other is labelled as Reason R.\\
Assertion A: Configuration auditing ensures that a change has been properly implemented.\\
Reason R: SCM is applied only after the software is delivered.
\optslong{Both A and R are true and R is the correct explanation of A}{Both A and R are true but R is NOT the correct explanation of A}{A is true but R is false}{A is false but R is true}
\end{question}

\begin{question}{M}
The SCM repository is
\optslong{a design pattern}{a test-case generator}{a compiler}{a database that stores configuration items, their versions and relationships}
\end{question}

\begin{question}{L}
Reverse engineering of data typically produces
\optslong{test cases}{data models such as entity-relationship diagrams from existing database structures}{a new user interface}{machine code}
\end{question}

\begin{question}{M}
Given below are two statements.\\
Statement I: Re-engineering is always cheaper than developing a new system from scratch.\\
Statement II: Inventory analysis in re-engineering assesses applications by business criticality, longevity and maintainability.
\optslong{Both Statement I and Statement II are true}{Both Statement I and Statement II are false}{Statement I is true but Statement II is false}{Statement I is false but Statement II is true}
\end{question}

\begin{answers}
\ans{1}{3}{Preventive maintenance is also called re-engineering.}
\ans{2}{2}{Studies typically attribute about half of maintenance effort to perfective changes.}
\ans{3}{1}{Baselines mark milestones in the delivery of SCIs.}
\ans{4}{1}{These are the SCM tasks described by Pressman.}
\ans{5}{3}{An approved change leads to an engineering change order (ECO).}
\ans{6}{2}{Data restructuring accompanies code restructuring; the model is cyclic.}
\ans{7}{3}{Reverse engineering is analysis only; it does not change the system.}
\ans{8}{4}{It often improves quality and adds new functions.}
\ans{9}{3}{SVN and CVS are centralised version control systems.}
\ans{10}{3}{Lehman's laws apply to E-type (real-world) systems.}
\ans{11}{2}{Estimates usually range from 60\% to 80\%.}
\ans{12}{3}{Completeness, interactivity and directionality are other properties of reverse engineering.}
\ans{13}{1}{Reuse also brings costs: finding, understanding and adapting components.}
\ans{14}{2}{COTS components are bought and integrated rather than built.}
\ans{15}{2}{Its three major activities are analysis, construction and dissemination.}
\ans{16}{4}{Legacy systems often have poor documentation, obsolete technology and degraded structure.}
\ans{17}{1}{It is similar to refactoring at a larger scale.}
\ans{18}{3}{Merging combines the changes from branches.}
\ans{19}{3}{SCM is an umbrella activity applied throughout the software process.}
\ans{20}{4}{It supports versioning, dependency tracking, requirements tracing and audit trails.}
\ans{21}{2}{It is often the first reverse-engineering task.}
\ans{22}{4}{Re-engineering is not always cheaper; inventory analysis helps decide which applications are worth re-engineering.}
\end{answers}
